\documentclass[a4paper,12pt]{article}
\usepackage[utf8]{inputenc}
\usepackage[polish]{babel}
\usepackage[T1]{fontenc}
\usepackage{amsmath}
\usepackage{graphicx}
\usepackage{geometry}
\usepackage{float}
\usepackage{caption}
\usepackage{listings}
\usepackage{xcolor}
\usepackage{subfigure}
\usepackage{amssymb}
\usepackage{natbib}

\geometry{margin=2.5cm}

\title{
\Large Zaawansowane metody przetwarzania i analizy sygnałów \\[6pt]
\textbf{Analiza częstotliwościowa HRV i klasyfikacja grup wiekowych }
}

\author{Piotr Sawicki 319003, Michał Odziemkowski 311267}
\date{Marzec 2026}

\begin{document}

\maketitle


\section{Cel projektu}
Głównym celem niniejszego projektu jest przeprowadzenie zaawansowanej analizy częstotliwościowej sygnału zmienności rytmu zatokowego (ang. \textit{Heart Rate Variability} – HRV)~. Projekt zakłada identyfikację istotnych statystycznie różnic w parametrach widmowych pomiędzy osobami młodymi a seniorami. Efektem końcowym ma być stworzenie stabilnego modelu klasyfikacyjnego, zdolnego do automatycznego rozróżniania grup wiekowych na podstawie cech wyekstrahowanych z sygnału EKG.

\section{Założenia projektowe}
W procesie badawczym przyjęto następujące założenia systemowe i merytoryczne:
\begin{itemize}
    \item Baza danych:
    
    Wykorzystanie zbioru danych \textit{Fantasia Database} (PhysioNet), zawierającego zbiór 120-minutowych zapisów EKG osób zdrowych w grupach: \textit{Young} (21–34 lata) oraz \textit{Elderly} (68–85 lat). Częstotliwość próbkowania sygnału wynosi $f_s = 250$ Hz.
    \item Dziedzina analizy:
    
    Skupienie się na parametrach częstotliwościowych (widmowych).
    \item Pasma częstotliwości:
    
    Analiza mocy w zakresach $LF$ (0.04--0.15 Hz) oraz $HF$ (0.15--0.4 Hz).
    \item Narzędzia:
    
    Implementacja algorytmów w środowisku MATLAB z wykorzystaniem dedykowanych bibliotek do przetwarzania sygnałów i uczenia maszynowego.
\end{itemize}

\newpage
\section{Metodyka i realizacja w środowisku MATLAB}
Realizacja projektu przebiega w pięciu kluczowych etapach:
\begin{enumerate}
    \item Ekstrakcja odstępów RR:
    
    Wstępne przetwarzanie sygnału EKG, filtracja zakłóceń oraz detekcja załamków R.
    %\item \textbf{Resampling i detrending:} Ponowne próbkowanie szeregu RR metodą splajnów (do częstotliwości 4 Hz) w celu uzyskania sygnału równomiernie spróbkowanego oraz usunięcie trendów wolnozmiennych.
    \item Estymacja widmowa: 
    
    Wyznaczenie gęstości widmowej mocy (PSD) przy użyciu metody Welcha (estymator nieparametryczny), co pozwala na redukcję wariancji szumu w porównaniu do klasycznej transformaty FFT.
    \item Wyznaczenie cech (Features): 
    
    Obliczenie wartości mocy pasm $LF$ i $HF$, całkowitej mocy widma (\textit{Total Power}) oraz wskaźnika balansu autonomicznego $LF/HF$.
    \item Klasyfikacja: 
    
    Zastosowanie algorytmu LDA (Linear Discriminant Analysis) lub SVM (Support Vector Machine) do podziału zbioru na klasy wiekowe.
\end{enumerate}

\section{Spodziewane wyniki i prezentacja danych}
W ramach końcowej fazy projektu przewiduje się uzyskanie następujących wyników:
\begin{itemize}
    \item Wizualizacja widma:
    
    Porównawcze wykresy gęstości widmowej mocy (PSD) dla reprezentatywnych pacjentów z obu grup, obrazujące spadek mocy składowej $HF$ u osób starszych.
    \item Analiza statystyczna:
    
    Zestawienie wartości parametrów $LF$ i $HF$ w formie wykresów.
    \item Ocena klasyfikatora:
    
    Wygenerowanie macierzy pomyłek (\textit{confusion matrix}), pozwalającej na ocenę czułości i swoistości modelu w rozpoznawaniu wieku pacjentów.
    \item Wnioski:
    
    Potwierdzenie hipotezy o spadku aktywności przywspółczulnej (składowa $HF$) wraz z wiekiem.
\end{itemize}

\end{document}